\chapter{Feature dictionary per dataset and cross-dataset mapping}
\label{app:features}
\chaptermark{Feature dictionary} %running head

This appendix lists the columns that every model receives, dataset by dataset, and the common
schema of the transfer of E9. A model input is every column of the data except the fine label,
the class and the columns named below for each dataset. The axioms read some of these columns
raw, before any coding. The numeric inputs are standardised and the categorical ones one-hot coded,
both fitted on the training part of the seed (\cref{sec:impl-data}). A category met only outside
the training part is coded as all zeros. The tables are written by
\texttt{analysis/feature\_dictionary.py} from the data and the knowledge bases of
\cref{app:axioms}.

\section{NSL-KDD}

KDDTrain+ and KDDTest+ \citep{tavallaee2009kddcup99} have 41 features, a label and a last
column with a difficulty level, which is not an input. The features are the basic
features of a connection, the content features of \cref{sec:content-host-based-features} and the
time-window and host-window aggregates of \cref{sec:time-window-connection-aggregates}.
\texttt{protocol\_type}, \texttt{service} and \texttt{flag} are categorical, and
\texttt{num\_outbound\_cmds} is zero in every record of KDDTrain+, so it carries nothing. After the
one-hot coding the network has 122 inputs, 121 for seed 2 (\cref{sec:impl-nslkdd}).
\Cref{tab:app-features-nslkdd} lists the columns.

\section{CIC-IDS2017}

The corrected version of \citet{liu2022error} is read as described in \cref{sec:impl-cic}. The
columns \texttt{id}, \texttt{Flow ID}, \texttt{Src IP}, \texttt{Src Port}, \texttt{Dst IP},
\texttt{Dst Port}, \texttt{Timestamp} and \texttt{Attempted Category} are not inputs: they name a
flow, its endpoints or its time, or they repeat the label. Without the destination port, a scan
whose flows differ only in the port scanned becomes a set of identical rows, which the cleaning
removes. \texttt{Protocol} is categorical, with the names of the protocol numbers. The 82 input
columns give the network 85 inputs (\cref{tab:app-features-cicids2017}).

\section{CIC-DDoS2019}

The columns \texttt{Flow ID}, \texttt{Source IP}, \texttt{Source Port}, \texttt{Destination IP},
\texttt{Destination Port}, \texttt{Timestamp} and the unnamed row index are not inputs, for the
same reasons, nor is \texttt{SimillarHTTP}, a text column. \texttt{Inbound} records the direction of the capture, not a
property of the traffic, and is kept in the data but never given to a model. The header length
written twice is read once (\cref{sec:impl-cic}). \texttt{Protocol} is categorical, and the 77
input columns give the network 79 inputs. Twelve columns hold one value in the whole training day:
the six bulk averages, three of the flag counts and three of the flag columns per direction
(\cref{tab:app-features-cicddos2019}). They carry nothing to the models trained on that day.

\section{Cross-dataset mapping}

The transfer of E9 trains on CIC-IDS2017 and tests on CIC-DDoS2019, whose flows come from two
versions of CICFlowMeter. The common schema of \cref{sec:impl-schema} matches the columns by a
normalised name and sixteen pairs by hand, and takes the names of CIC-DDoS2019.
\Cref{tab:app-transfer-schema} lists the nineteen columns of CIC-IDS2017 that change name. Every
column of CIC-DDoS2019 has a counterpart, and five columns of CIC-IDS2017 have none and stay out
of the transfer: \texttt{Fwd RST Flags}, \texttt{Bwd RST Flags}, \texttt{ICMP Code},
\texttt{ICMP Type} and \texttt{Total TCP Flow Time}. The common schema has 77 columns. Equal names do not
make equal measures: the extractors differ, so a column with one name can be computed differently
on the two sides.

% written by analysis/feature_dictionary.py, do not edit by hand

\begin{table}[p]
\centering
\scriptsize
\caption{NSL-KDD: input columns in the order of the data. cat., $k$: categorical with $k$ values in KDDTrain+; const.: one value in all of it; axiom identifiers: axioms of the main knowledge base that read the column (\cref{app:axioms}).}
\label{tab:app-features-nslkdd}
\setlength{\tabcolsep}{3pt}
\begin{tabular}{@{}llll@{}}
\toprule
Feature & Note & Feature & Note \\
\midrule
\texttt{duration} &  & \texttt{is\_guest\_login} &  \\
\texttt{protocol\_type} & cat., 3 & \texttt{count} & A1 \\
\texttt{service} & cat., 70 & \texttt{srv\_count} &  \\
\texttt{flag} & cat., 11; A1, A3 & \texttt{serror\_rate} &  \\
\texttt{src\_bytes} &  & \texttt{srv\_serror\_rate} &  \\
\texttt{dst\_bytes} &  & \texttt{rerror\_rate} &  \\
\texttt{land} &  & \texttt{srv\_rerror\_rate} &  \\
\texttt{wrong\_fragment} &  & \texttt{same\_srv\_rate} &  \\
\texttt{urgent} &  & \texttt{diff\_srv\_rate} &  \\
\texttt{hot} &  & \texttt{srv\_diff\_host\_rate} &  \\
\texttt{num\_failed\_logins} &  & \texttt{dst\_host\_count} &  \\
\texttt{logged\_in} &  & \texttt{dst\_host\_srv\_count} &  \\
\texttt{num\_compromised} &  & \texttt{dst\_host\_same\_srv\_rate} &  \\
\texttt{root\_shell} &  & \texttt{dst\_host\_diff\_srv\_rate} & A1, A2 \\
\texttt{su\_attempted} &  & \texttt{dst\_host\_same\_src\_port\_rate} &  \\
\texttt{num\_root} &  & \texttt{dst\_host\_srv\_diff\_host\_rate} &  \\
\texttt{num\_file\_creations} &  & \texttt{dst\_host\_serror\_rate} &  \\
\texttt{num\_shells} &  & \texttt{dst\_host\_srv\_serror\_rate} &  \\
\texttt{num\_access\_files} &  & \texttt{dst\_host\_rerror\_rate} &  \\
\texttt{num\_outbound\_cmds} & const. & \texttt{dst\_host\_srv\_rerror\_rate} &  \\
\texttt{is\_host\_login} &  &  &  \\
\bottomrule
\end{tabular}
\end{table}

\begin{table}[p]
\centering
\scriptsize
\caption{CIC-IDS2017: input columns in the order of the data. cat., $k$: categorical with $k$ values in the whole dataset; const.: one value in all of it; axiom identifiers: axioms of the main knowledge base that read the column (\cref{app:axioms}).}
\label{tab:app-features-cicids2017}
\setlength{\tabcolsep}{3pt}
\begin{tabular}{@{}llll@{}}
\toprule
Feature & Note & Feature & Note \\
\midrule
\texttt{Protocol} & cat., 4 & \texttt{Packet Length Max} &  \\
\texttt{Flow Duration} &  & \texttt{Packet Length Mean} &  \\
\texttt{Total Fwd Packet} &  & \texttt{Packet Length Std} &  \\
\texttt{Total Bwd packets} &  & \texttt{Packet Length Variance} &  \\
\texttt{Total Length of Fwd Packet} & A1 & \texttt{FIN Flag Count} &  \\
\texttt{Total Length of Bwd Packet} &  & \texttt{SYN Flag Count} &  \\
\texttt{Fwd Packet Length Max} & A1, A2 & \texttt{RST Flag Count} &  \\
\texttt{Fwd Packet Length Min} &  & \texttt{PSH Flag Count} &  \\
\texttt{Fwd Packet Length Mean} &  & \texttt{ACK Flag Count} &  \\
\texttt{Fwd Packet Length Std} &  & \texttt{URG Flag Count} &  \\
\texttt{Bwd Packet Length Max} &  & \texttt{CWR Flag Count} &  \\
\texttt{Bwd Packet Length Min} &  & \texttt{ECE Flag Count} &  \\
\texttt{Bwd Packet Length Mean} &  & \texttt{Down/Up Ratio} &  \\
\texttt{Bwd Packet Length Std} & A2 & \texttt{Average Packet Size} &  \\
\texttt{Flow Bytes/s} &  & \texttt{Fwd Segment Size Avg} &  \\
\texttt{Flow Packets/s} &  & \texttt{Bwd Segment Size Avg} & A1 \\
\texttt{Flow IAT Mean} &  & \texttt{Fwd Bytes/Bulk Avg} &  \\
\texttt{Flow IAT Std} &  & \texttt{Fwd Packet/Bulk Avg} &  \\
\texttt{Flow IAT Max} &  & \texttt{Fwd Bulk Rate Avg} &  \\
\texttt{Flow IAT Min} &  & \texttt{Bwd Bytes/Bulk Avg} &  \\
\texttt{Fwd IAT Total} &  & \texttt{Bwd Packet/Bulk Avg} &  \\
\texttt{Fwd IAT Mean} &  & \texttt{Bwd Bulk Rate Avg} &  \\
\texttt{Fwd IAT Std} &  & \texttt{Subflow Fwd Packets} &  \\
\texttt{Fwd IAT Max} &  & \texttt{Subflow Fwd Bytes} &  \\
\texttt{Fwd IAT Min} &  & \texttt{Subflow Bwd Packets} &  \\
\texttt{Bwd IAT Total} &  & \texttt{Subflow Bwd Bytes} &  \\
\texttt{Bwd IAT Mean} &  & \texttt{FWD Init Win Bytes} &  \\
\texttt{Bwd IAT Std} &  & \texttt{Bwd Init Win Bytes} &  \\
\texttt{Bwd IAT Max} &  & \texttt{Fwd Act Data Pkts} &  \\
\texttt{Bwd IAT Min} &  & \texttt{Fwd Seg Size Min} &  \\
\texttt{Fwd PSH Flags} &  & \texttt{Active Mean} &  \\
\texttt{Bwd PSH Flags} & A2 & \texttt{Active Std} &  \\
\texttt{Fwd URG Flags} &  & \texttt{Active Max} &  \\
\texttt{Bwd URG Flags} &  & \texttt{Active Min} &  \\
\texttt{Fwd RST Flags} &  & \texttt{Idle Mean} &  \\
\texttt{Bwd RST Flags} &  & \texttt{Idle Std} &  \\
\texttt{Fwd Header Length} &  & \texttt{Idle Max} &  \\
\texttt{Bwd Header Length} &  & \texttt{Idle Min} &  \\
\texttt{Fwd Packets/s} &  & \texttt{ICMP Code} &  \\
\texttt{Bwd Packets/s} &  & \texttt{ICMP Type} &  \\
\texttt{Packet Length Min} &  & \texttt{Total TCP Flow Time} &  \\
\bottomrule
\end{tabular}
\end{table}

\begin{table}[p]
\centering
\scriptsize
\caption{CIC-DDoS2019: input columns in the order of the data. cat., $k$: categorical with $k$ values in the sampled training day; const.: one value in all of it; axiom identifiers: axioms of the main knowledge base that read the column (\cref{app:axioms}).}
\label{tab:app-features-cicddos2019}
\setlength{\tabcolsep}{3pt}
\begin{tabular}{@{}llll@{}}
\toprule
Feature & Note & Feature & Note \\
\midrule
\texttt{Protocol} & cat., 3; A1 & \texttt{Max Packet Length} &  \\
\texttt{Flow Duration} &  & \texttt{Packet Length Mean} &  \\
\texttt{Total Fwd Packets} &  & \texttt{Packet Length Std} &  \\
\texttt{Total Backward Packets} & A1 & \texttt{Packet Length Variance} &  \\
\texttt{Total Length of Fwd Packets} &  & \texttt{FIN Flag Count} & const. \\
\texttt{Total Length of Bwd Packets} &  & \texttt{SYN Flag Count} &  \\
\texttt{Fwd Packet Length Max} &  & \texttt{RST Flag Count} &  \\
\texttt{Fwd Packet Length Min} &  & \texttt{PSH Flag Count} & const. \\
\texttt{Fwd Packet Length Mean} &  & \texttt{ACK Flag Count} &  \\
\texttt{Fwd Packet Length Std} &  & \texttt{URG Flag Count} &  \\
\texttt{Bwd Packet Length Max} &  & \texttt{CWE Flag Count} &  \\
\texttt{Bwd Packet Length Min} &  & \texttt{ECE Flag Count} & const. \\
\texttt{Bwd Packet Length Mean} &  & \texttt{Down/Up Ratio} &  \\
\texttt{Bwd Packet Length Std} &  & \texttt{Average Packet Size} &  \\
\texttt{Flow Bytes/s} &  & \texttt{Avg Fwd Segment Size} &  \\
\texttt{Flow Packets/s} & A2 & \texttt{Avg Bwd Segment Size} &  \\
\texttt{Flow IAT Mean} &  & \texttt{Fwd Avg Bytes/Bulk} & const. \\
\texttt{Flow IAT Std} &  & \texttt{Fwd Avg Packets/Bulk} & const. \\
\texttt{Flow IAT Max} &  & \texttt{Fwd Avg Bulk Rate} & const. \\
\texttt{Flow IAT Min} &  & \texttt{Bwd Avg Bytes/Bulk} & const. \\
\texttt{Fwd IAT Total} &  & \texttt{Bwd Avg Packets/Bulk} & const. \\
\texttt{Fwd IAT Mean} &  & \texttt{Bwd Avg Bulk Rate} & const. \\
\texttt{Fwd IAT Std} &  & \texttt{Subflow Fwd Packets} &  \\
\texttt{Fwd IAT Max} &  & \texttt{Subflow Fwd Bytes} &  \\
\texttt{Fwd IAT Min} &  & \texttt{Subflow Bwd Packets} &  \\
\texttt{Bwd IAT Total} &  & \texttt{Subflow Bwd Bytes} &  \\
\texttt{Bwd IAT Mean} &  & \texttt{Init\_Win\_bytes\_forward} &  \\
\texttt{Bwd IAT Std} &  & \texttt{Init\_Win\_bytes\_backward} &  \\
\texttt{Bwd IAT Max} &  & \texttt{act\_data\_pkt\_fwd} &  \\
\texttt{Bwd IAT Min} &  & \texttt{min\_seg\_size\_forward} &  \\
\texttt{Fwd PSH Flags} &  & \texttt{Active Mean} &  \\
\texttt{Bwd PSH Flags} & const. & \texttt{Active Std} &  \\
\texttt{Fwd URG Flags} & const. & \texttt{Active Max} &  \\
\texttt{Bwd URG Flags} & const. & \texttt{Active Min} &  \\
\texttt{Fwd Header Length} &  & \texttt{Idle Mean} &  \\
\texttt{Bwd Header Length} &  & \texttt{Idle Std} &  \\
\texttt{Fwd Packets/s} &  & \texttt{Idle Max} &  \\
\texttt{Bwd Packets/s} &  & \texttt{Idle Min} &  \\
\texttt{Min Packet Length} &  &  &  \\
\bottomrule
\end{tabular}
\end{table}

\begin{table}[p]
\centering
\scriptsize
\caption{Common schema of the transfer: columns of CIC-IDS2017 whose name changes, with the name of CIC-DDoS2019 they take. 77 columns are common.}
\label{tab:app-transfer-schema}
\setlength{\tabcolsep}{4pt}
\begin{tabular}{@{}ll@{}}
\toprule
CIC-IDS2017 & CIC-DDoS2019 \\
\midrule
\texttt{Total Fwd Packet} & \texttt{Total Fwd Packets} \\
\texttt{Total Bwd packets} & \texttt{Total Backward Packets} \\
\texttt{Total Length of Fwd Packet} & \texttt{Total Length of Fwd Packets} \\
\texttt{Total Length of Bwd Packet} & \texttt{Total Length of Bwd Packets} \\
\texttt{Packet Length Min} & \texttt{Min Packet Length} \\
\texttt{Packet Length Max} & \texttt{Max Packet Length} \\
\texttt{CWR Flag Count} & \texttt{CWE Flag Count} \\
\texttt{Fwd Segment Size Avg} & \texttt{Avg Fwd Segment Size} \\
\texttt{Bwd Segment Size Avg} & \texttt{Avg Bwd Segment Size} \\
\texttt{Fwd Bytes/Bulk Avg} & \texttt{Fwd Avg Bytes/Bulk} \\
\texttt{Fwd Packet/Bulk Avg} & \texttt{Fwd Avg Packets/Bulk} \\
\texttt{Fwd Bulk Rate Avg} & \texttt{Fwd Avg Bulk Rate} \\
\texttt{Bwd Bytes/Bulk Avg} & \texttt{Bwd Avg Bytes/Bulk} \\
\texttt{Bwd Packet/Bulk Avg} & \texttt{Bwd Avg Packets/Bulk} \\
\texttt{Bwd Bulk Rate Avg} & \texttt{Bwd Avg Bulk Rate} \\
\texttt{FWD Init Win Bytes} & \texttt{Init\_Win\_bytes\_forward} \\
\texttt{Bwd Init Win Bytes} & \texttt{Init\_Win\_bytes\_backward} \\
\texttt{Fwd Act Data Pkts} & \texttt{act\_data\_pkt\_fwd} \\
\texttt{Fwd Seg Size Min} & \texttt{min\_seg\_size\_forward} \\
\bottomrule
\end{tabular}
\end{table}

% only CIC-IDS2017: Fwd RST Flags, Bwd RST Flags, ICMP Code, ICMP Type, Total TCP Flow Time
% only CIC-DDoS2019: none

